% 卡齐米日·库拉托夫斯基（Kazimierz Kuratowski）（综述）
% license CCBYSA3
% type Wiki

本文根据 CC-BY-SA 协议转载翻译自维基百科\href{https://en.wikipedia.org/wiki/Kazimierz_Kuratowski}{相关文章}。

\begin{figure}[ht]
\centering
\includegraphics[width=6cm]{./figures/06e709d21f762e7a.png}
\caption{1959年的库拉托夫斯基} \label{fig_KQKLS_2}
\end{figure}
卡齐米日·库拉托夫斯基（Kazimierz Kuratowski，波兰语发音：[kaˈʑimjɛʂ kuraˈtɔfskʲi]；1896年2月2日－1980年6月18日）是一位波兰数学家和逻辑学家。他是华沙数学学派的主要代表人物之一，曾任华沙大学教授以及波兰科学院数学研究所（IM PAN）的研究员。1946年至1953年间，他担任波兰数学学会会长。

他以在集合论、拓扑学、测度论和图论方面的贡献而著称。以他名字命名的一些著名数学概念包括：库拉托夫斯基定理、库拉托夫斯基闭包公理、库拉托夫斯基–佐恩引理以及库拉托夫斯基交集定理。
\subsection{生平与事业}
\subsubsection{早年生活}
卡齐米日·库拉托夫斯基于1896年2月2日出生在华沙（当时属于受俄罗斯帝国控制的波兰会议王国）。他的父亲马雷克·库拉托夫是一名律师，母亲是罗莎·卡热夫斯卡。他在以保卫将军帕韦乌·赫沙诺夫斯基（Paweł Chrzanowski）命名的华沙中学完成中等教育。1913年，他进入苏格兰的格拉斯哥大学学习工程学，部分原因是他不愿在俄语教学体系下学习——当时在波兰禁止使用波兰语授课【需要引用】。他仅学习了一年，第一次世界大战的爆发使进一步的学业中断。1915年，俄军撤出华沙，华沙大学重新开学，并以波兰语作为教学语言。同年，库拉托夫斯基重返大学，这次主修数学。1921年，他在新成立的第二波兰共和国获得博士学位。
\subsubsection{博士论文}
1921年秋，库拉托夫斯基因其开创性的研究工作获得博士学位。他的论文由两部分组成。第一部分致力于通过闭包公理建立拓扑学的公理化体系。这一部分（1922年以略作修改的形式再版）被数百篇科学论文引用过\(^\text{[1]}\)。

库拉托夫斯基论文的第二部分研究的是“两点间不可约连续体”的问题。这一主题最初是齐格蒙特·亚尼谢夫斯基（Zygmunt Janiszewski）在其法语博士论文中的研究对象。由于亚尼谢夫斯基已去世，库拉托夫斯基的论文导师改由斯特凡·马祖凯维奇（Stefan Mazurkiewicz）担任。库拉托夫斯基的研究解决了比利时数学家夏尔-让·艾蒂安·古斯塔夫·尼古拉·德·拉瓦莱-普桑男爵（Charles-Jean Étienne Gustave Nicolas, Baron de la Vallée Poussin）在集合论中提出的若干问题。
\subsubsection{二战前的学术生涯}
两年后的1923年，库拉托夫斯基被任命为华沙大学数学系副教授。1927年，他受聘为利沃夫理工大学（Lwów Polytechnic）的数学正教授，并一直担任该校数学系主任至1933年。库拉托夫斯基还曾两度出任该系的院长。1929年，他成为华沙科学学会的成员。

虽然库拉托夫斯基与利沃夫数学学派的许多学者，如斯特凡·巴拿赫（Stefan Banach）和斯坦尼斯瓦夫·乌拉姆（Stanisław Ulam），以及以苏格兰咖啡馆（Scottish Café）为中心的一群数学家有着密切的学术往来，但他始终与华沙保持紧密联系。1934年，库拉托夫斯基离开利沃夫返回华沙，那是在著名的《苏格兰问题集》（Scottish Book）开始编写（1935年）之前，因此他并未向其中贡献任何问题。然而，他确实与巴拿赫在测度论的重要问题上有过密切合作\(^\text{[2][3]}\)。

1934年，他被任命为华沙大学教授。一年后，库拉托夫斯基被任命为数学系主任。1936年至1939年间，他担任波兰科学与应用科学委员会数学委员会秘书。
\subsubsection{战时与战后}
第二次世界大战期间，由于德国占领当局禁止波兰人接受高等教育，库拉托夫斯基在华沙的地下大学秘密授课。

1945年2月，华沙大学重新开学后，库拉托夫斯基开始在该校讲授课程。同年，他成为波兰学习科学院（Polish Academy of Learning）院士。1946年，他被任命为华沙大学数学系副主任，1949年起，他当选为华沙科学学会副会长。1952年，他成为波兰科学院（Polish Academy of Sciences）院士，并在1957年至1968年间担任该院副院长。

二战后，库拉托夫斯基积极参与波兰科学界的重建。他协助创立了国家数学研究所（State Institute of Mathematics），该机构于1952年并入波兰科学院\(^\text{[4]}\)。1948年至1967年间，库拉托夫斯基担任波兰科学院数学研究所所长，同时长期担任波兰数学学会及国际数学学会的主席。他还在1963年至1966年间担任国际数学联盟（International Mathematical Union）副主席，并于1968年至1980年出任国家数学研究所科学委员会主席。1948年至1980年间，他一直领导该院的拓扑学部门。

他的学生之一是著名逻辑学家兼集合论学者安杰伊·莫斯托夫斯基（Andrzej Mostowski）。
\subsection{遗产与影响}
卡齐米日·库拉托夫斯基是著名的波兰数学家群体成员之一，他们常在利沃夫的“苏格兰咖啡馆”（Scottish Café）聚会讨论数学问题。他曾任波兰数学学会（Polish Mathematical Society, PTM）会长，也是华沙科学学会（Warsaw Scientific Society, TNW）的成员。此外，他还担任《数学基础》（Fundamenta Mathematicae）的主编——该刊属于《波兰数学学会年刊》（Annals of the Polish Mathematical Society）系列出版物。与此同时，库拉托夫斯基还担任《波兰科学院公报》（Bulletin of the Polish Academy of Sciences）的编辑之一。他也是《数学专著》（Mathematical Monographs）系列丛书的主要撰稿人之一，该系列由波兰科学院数学研究所（Institute of Mathematics of the Polish Academy of Sciences, IMPAN）合作出版。该丛书收录了来自华沙学派和利沃夫学派的高质量研究专著，内容涵盖纯数学与应用数学的各个领域。

卡齐米日·库拉托夫斯基还是多个科学学会及国外科学院的活跃成员，其中包括爱丁堡皇家学会，以及奥地利、德国、匈牙利、意大利和苏联（USSR）的学术机构。
\subsubsection{卡齐米日·库拉托夫斯基奖}
1981年，波兰科学院数学研究所（IMPAN）、波兰数学学会以及库拉托夫斯基的女儿佐菲娅·库拉托夫斯卡（Zofia Kuratowska）共同设立了以他名字命名的奖项——库拉托夫斯基奖，用于表彰30岁以下在数学领域取得杰出成就的青年学者\(^\text{[5]}\)。该奖被认为是授予年轻波兰数学家的最高荣誉之一。历届获奖者包括约瑟夫·H·普日蒂茨基（Józef H. Przytycki）、马留什·莱曼奇克（Mariusz Lemańczyk）、托马什·武恰克（Tomasz Łuczak）、米科瓦伊·博扬奇克（Mikołaj Bojańczyk）以及沃伊切赫·萨莫季（Wojciech Samotij）等\(^\text{[5]}\)。
\subsection{研究}
\begin{figure}[ht]
\centering
\includegraphics[width=6cm]{./figures/b0b5289ce14cc580.png}
\caption{无需文字的证明：利用库拉托夫斯基定理或瓦格纳定理，通过在超立方体图中找到 K₅（上图）或 K₃,₃（下图）子图，证明其为非平面图。} \label{fig_KQKLS_1}
\end{figure}
库拉托夫斯基的研究主要集中于抽象拓扑与度量结构。他提出并形式化了闭包公理（在数学界被称为库拉托夫斯基闭包公理），这一体系为拓扑空间理论及“两点间不可约连续体理论”的发展奠定了基础。战后，库拉托夫斯基最重要的研究成果之一是关于拓扑学与解析函数（理论）之间关系的研究，以及欧几里得空间切割问题领域的探索。

在利沃夫时期，他最杰出的学生之一乌拉姆（Ulam）与他合作，引入了所谓的拟同胚（quasi-homeomorphism）概念，为拓扑学研究开辟了一个全新的方向。库拉托夫斯基在测度论方面的研究（包括与巴拿赫和塔斯基的合作）被许多学生进一步发展。此外，他与阿尔弗雷德·塔斯基（Alfred Tarski）及瓦茨瓦夫·谢尔宾斯基（Wacław Sierpiński）共同构建了关于波兰空间（Polish spaces）的大部分理论——这些空间的名称正是为了纪念他们的贡献与学术遗产。克纳斯特（Knaster）与库拉托夫斯基还对连通分支理论进行了系统而精确的研究，并将其应用于诸如平面切割问题等领域，其中包括一些连通分支的悖论性实例。

1922年，库拉托夫斯基证明了后来被称为库拉托夫斯基–佐恩引理（Kuratowski–Zorn Lemma）的结果（通常简称为佐恩引理）\(^\text{[6]}\)。这一结果与许多基础数学定理密切相关。1935年，佐恩（Zorn）给出了该引理的应用\(^\text{[7]}\)。库拉托夫斯基在集合论与拓扑学中引入了许多重要概念，并在多个领域建立了新的术语与符号体系。他对数学的主要贡献包括：
\begin{itemize}
\item 对拓扑空间的刻画，即现在所称的库拉托夫斯基闭包公理；
\item 库拉托夫斯基–佐恩引理的证明；
\item 在图论中，对平面图的刻画，即如今被称为库拉托夫斯基定理的结果；
\item 将有序对$x, y$表示为集合${{x}, {x, y}}$的定义；\(^\text{[8]}\)
\item 库拉托夫斯基有限集的定义，参见 Kuratowski-finite；\(^\text{[9]}\)
\item 塔斯基–库拉托夫斯基算法的提出；
\item 库拉托夫斯基的闭包–补集问题；
\item 库拉托夫斯基嵌入；
\item 库拉托夫斯基自由集定理；
\item 库拉托夫斯基交集定理；
\item 克纳斯特–库拉托夫斯基扇形（fan）；
\item 库拉托夫斯基–乌拉姆定理；
\item 度量空间子集的库拉托夫斯基收敛；
\item 库拉托夫斯基与雷尔–纳尔泽夫斯基可测选择定理。
\end{itemize}
库拉托夫斯基战后的研究主要集中在三个方向：
\begin{itemize}
\item 连续函数中的同伦理论的发展；
\item 高维空间中连通空间理论的构建；
\item 基于集合连续变换性质，对欧几里得空间被其任意子集切割的统一描述。
\end{itemize}
\subsection{出版著作}
库拉托夫斯基发表的170多篇作品中，包括多部重要的专著与书籍，如《拓扑学》（Topologie，第一卷1933年出版，被译为英文和俄文；第二卷1950年出版）以及《集合论与拓扑学导论》（Introduction to Set Theory and Topology，第一卷1952年出版，被译为英文、法文、西班牙文和保加利亚文）。他还著有《波兰数学的半个世纪：1920–1970 回忆与反思》（A Half Century of Polish Mathematics 1920–1970: Remembrances and Reflections，1973年）\(^\text{[10]}\)与《自传笔记》（Notes to his autobiography，1981年）。后者在他去世后由其女儿佐菲娅·库拉托夫斯卡（Zofia Kuratowska）整理出版。

卡齐米日·库拉托夫斯基曾代表波兰数学界参与国际数学联盟（International Mathematical Union）的工作，并在1963年至1966年间担任副主席。此外，他还多次参加国际学术大会，并在世界各地数十所大学讲学。他被授予格拉斯哥大学、布拉格大学、弗罗茨瓦夫大学及巴黎大学的名誉博士学位，并获得波兰最高国家荣誉、捷克斯洛伐克科学院金质奖章以及波兰科学院金质奖章。库拉托夫斯基于1980年6月18日逝世于华沙。

库拉托夫斯基，卡齐米日；莫斯托夫斯基，安杰伊（1976）[1968]，《集合论：含描述集合论导论》，载于《逻辑与数学基础研究》（Studies in Logic and the Foundations of Mathematics），第86卷（第二版），阿姆斯特丹–纽约–牛津：北荷兰出版社（North-Holland Publishing Co.），MR 0485384。
\subsection{参见}
\begin{itemize}
\item 波兰数学家名单
\item 苏格兰咖啡馆
\item 以卡齐米日·库拉托夫斯基命名的事物列表
\item 波兰科学与技术发展时间线
\end{itemize}
\subsection{注释}
\begin{enumerate}
\item Kuratowski 1922.
\item MacTutor 文章：Kazimierz Kuratowski.
\item 《[www.day.kyiv.ua》文章：“苏格兰笔记本：利沃夫的数学遗产”。](http://www.day.kyiv.ua》文章：“苏格兰笔记本：利沃夫的数学遗产”。)
\item “历史”。impan.pl。存档于2024年12月25日，访问于2024年12月7日。
\item “卡齐米日·库拉托夫斯基奖”。波兰科学院数学研究所（波兰语）。存档于2020年11月7日，访问于2017年7月19日。
\item Fundamenta Mathematicae 第3卷，第页。
\item Bulletin of the American Mathematical Society，第41卷，第页。
\item Kuratowski 1921，第171页。
\item Kuratowski 1920。
\item Kuratowski 1980。
\end{enumerate}
\subsection{参考文献}
\begin{itemize}
\item Borsuk, Karol (1960)，《论卡齐米日·库拉托夫斯基教授在拓扑学方面的成就》（波兰语），《数学通讯》（iadomości Matematyczne，第2卷第3期，页231–237。
\item Kuratowski, Kazimierz (1920)，《关于有限集的概念》（Sur la notion d'ensemble fini），Fundamenta Mathematicae，第1卷，页129–131。doi:10.4064/fm-1-1-129-131。
\item Kuratowski, Kazimierz (1921)，《关于集合论中的序概念》（Sur la notion de l'ordre dans la Théorie des Ensembles），Fundamenta Mathematicae*，第2卷第1期，页161–171。doi:10.4064/fm-2-1-161-171。
\item Kuratowski, Kazimierz (1922)，《关于解析位置论中的Ā运算》（Sur l'opération Ā de l'Analysis Situs），Fundamenta Mathematicae，第3卷，页182–199。doi:10.4064/fm-3-1-182-199。ISSN 0016-2736。
\item Kuratowski, Kazimierz (1980)，《波兰数学的半个世纪：回忆与反思》，牛津：Pergamon出版社。ISBN 978-0-08-023046-7。
\end{itemize}
\subsection{外部链接}
\begin{itemize}
\item 《拓扑学 I：可度量空间与完备空间》（TOPOLOGIE I, Espaces Métrisables, Espaces Complets），《数学专著》（Monografie Matematyczne）系列，第20卷，波兰数学学会，华沙–利沃夫，1948年。
\item 《拓扑学 II：紧致空间、连通空间与欧几里得平面》（TOPOLOGIE II, Espaces Compacts, Espaces Connexes, Plan Euclidien），《数学专著》系列，第21卷，波兰数学学会，华沙–利沃夫，1950年。
\item O’Connor, John J.; Robertson, Edmund F., “Kazimierz Kuratowski”，圣安德鲁斯大学《MacTutor 数学史档案》。
\item 数学家族谱项目（Mathematics Genealogy Project）中的卡齐米日·库拉托夫斯基。
\end{itemize}